% Lecture example set: SC4061-L01
% Source: L1 Imaging concepts; questions are self-contained check examples
% Human verified: Yes

\exercisemeta
  {SC4061-L01-EX}
  {2026-08-11}
  {L1 Imaging 的 thin lens、sampling 与 quantization 内容；题目为自编检验例}
  {答案已完成并确认}
  {已确认（2026-08-11）}

\exercisequestion{SC4061-L01-E01}{Thin Lens 成像与放大率}

\textbf{对应知识点：}
\relatedtopic{SC4061-L01-T03}{Thin Lens Model 与 Magnification}

\subsubsection*{题目}

一个 focal length（焦距）为 $f=50\,\mathrm{mm}$ 的 thin lens（薄透镜）拍摄距离 $u=500\,\mathrm{mm}$、高度 $Y=180\,\mathrm{mm}$ 的物体。求 image distance（像距）$v$、signed magnification（带符号放大率）$m$ 和像高 $y$。

\begin{workedsolution}
由 thin lens equation（薄透镜公式）
\[
  \frac{1}{v}=\frac{1}{f}-\frac{1}{u}
  =\frac{1}{50}-\frac{1}{500}
  =\frac{9}{500},
\]
得到
\[
  v=\frac{500}{9}\,\mathrm{mm}\approx55.6\,\mathrm{mm}.
\]
因此
\[
  m=-\frac{v}{u}=-\frac{1}{9}\approx-0.111,
  \qquad
  y=mY=-20\,\mathrm{mm}.
\]
Sensor（传感器）应位于透镜后约 $55.6\,\mathrm{mm}$；像高为 $20\,\mathrm{mm}$，负号表示 image（像）相对物体倒立。
\end{workedsolution}

\textbf{检查点：}只写 $|m|$ 会丢失 orientation（方向）信息；采用带符号约定时，real inverted image（倒立实像）的放大率为负。

\exercisequestion{SC4061-L01-E02}{Sampling、Quantization 与存储量}

\textbf{对应知识点：}
\relatedtopic{SC4061-L01-T06}{Sensor、Sampling、Quantization 与 Image Matrix}

\subsubsection*{题目}

一幅未压缩 RGB image（RGB 彩色图像）包含 $1{,}000{,}000$ 个 pixels（像素），每个 channel（通道）采用 8-bit quantization（8 位量化）。

\begin{enumerate}[label=(\alph*)]
  \item 求原始存储量，并分别用 MB（十进制）和 MiB（二进制）表示。
  \item 若两个空间方向的 sampling interval（采样间隔）都减半，而 field of view（视野）不变，pixel count（像素数）变为原来的多少倍？
  \item 分别写出 sampling rate（采样率）不足和 quantization levels（量化等级）不足的典型失真。
\end{enumerate}

\begin{workedsolution}
每个 pixel 包含三个 8-bit channels，因此
\[
  S=1{,}000{,}000\times3\times8
   =24{,}000{,}000\ \text{bits}
   =3{,}000{,}000\ \text{bytes}.
\]
所以存储量为 $3.00\,\mathrm{MB}$，或
\[
  \frac{3{,}000{,}000}{2^{20}}\approx2.86\,\mathrm{MiB}.
\]
两个方向的间隔都减半时，每个方向的采样点数加倍，总像素数变为 $2\times2=4$ 倍。Sampling rate 不足会产生 aliasing（混叠）；quantization levels 不足会产生 banding（色带）或明显的灰度跳变。
\end{workedsolution}

\textbf{检查点：}“1 megapixel”本身不能确定文件大小；还必须说明 channel count（通道数）、bit depth（位深）以及是否 compression（压缩）。
